% Requires: \usepackage{amsmath,amssymb}
\section{Pareto Frontier Scoring}
\label{sec:pareto-scoring}

Each evaluated design is optimized by the agent over $T$ iterations. Iteration
$i$ submits a small set of parameterized implementations, and we retain the set
$S_i$ of \emph{scoreable} points, i.e., those that render, pass the unchanged
testbench, and report synthesis metrics. The unmodified design is measured once
as a baseline $b$. Every point $p$ has a worst-case latency $L(p)$ in cycles and
a usage $R_r(p)$ of each resource type $r \in \{\text{LUT}, \text{FF}, \text{DSP},
\text{BRAM}\}$. Because different designs have very different scales and
tradeoff curves, we want a single scalar that says how good the tradeoff
frontier found by the agent is, that is comparable across designs, and that
credits both the fast/large and the slow/small ends of the frontier.
We obtain it from the hypervolume of the latency--resource Pareto front.

\paragraph{Objective space.}
We score each resource type separately, in the two-dimensional space
$(R_r, L)$ where both objectives are minimized. All values are mapped to
$\tilde{x} = \log_{10}(x + 1)$ before any front or area is computed. The
logarithm makes improvements count by ratio rather than by absolute amount
(halving latency is worth the same at $10^3$ and at $10^6$ cycles), and the
$+1$ keeps a resource count of zero well defined. A point $p$ \emph{dominates}
$q$ if $\tilde{R}_r(p) \le \tilde{R}_r(q)$ and $\tilde{L}(p) \le \tilde{L}(q)$
with at least one strict inequality. For a set of points $P$, we write
$\mathrm{ND}(P)$ for its Pareto front, the subset of points not dominated by
any other point in $P$.

\paragraph{Hypervolume.}
For a front $F$ and a reference point $\rho = (\rho_R, \rho_L)$, the
hypervolume $\mathrm{HV}(F)$ is the area of the region that is dominated by $F$
and bounded by $\rho$. If $F$ is sorted by increasing resource,
$\tilde{R}_1 < \dots < \tilde{R}_k$, with strictly decreasing latency
$\tilde{L}_1 > \dots > \tilde{L}_k$, then
\begin{equation}
  \mathrm{HV}(F) = \sum_{j=1}^{k} \bigl(\rho_R - \tilde{R}_j\bigr)
  \bigl(\tilde{L}_{j-1} - \tilde{L}_j\bigr), \qquad \tilde{L}_0 \equiv \rho_L .
\end{equation}
A larger hypervolume means the front reaches lower latency, lower resource use,
or covers more of the tradeoff between them.

\paragraph{Reference point and normalization.}
For each design and resource type, let
\begin{equation}
  P^{*} = \mathrm{ND}\Bigl(\{b\} \cup \bigcup_{i=1}^{T} S_i\Bigr)
\end{equation}
be the combined front of the baseline and \emph{all} points found in
\emph{all} iterations. The reference point $\rho$ is the worst latency and worst
resource use among points of $P^{*}$, each increased by a $1.1\times$ margin
on $x+1$ (an additive $\log_{10} 1.1$ in log space), so that the extreme points
of $P^{*}$ contribute a nonzero area. The reference is fixed per design and
resource, which makes scores of different iterations of the same design
directly comparable. The score of a point set $P$ is
\begin{equation}
  s(P) \;=\; \frac{\mathrm{HV}\bigl(\mathrm{ND}(P \cup \{b\})\bigr) - \mathrm{HV}(\{b\})}
                 {\mathrm{HV}(P^{*}) - \mathrm{HV}(\{b\})} .
  \label{eq:score}
\end{equation}
The score is anchored at the baseline: $s = 0$ means $P$ adds nothing to the
baseline, and $s = 1$ means $P$ recovers the entire combined front $P^{*}$. For
any $P \subseteq \bigcup_i S_i$ it lies in $[0, 1]$. It therefore reports the
fraction of the improvement over the baseline that was actually achieved, and
it is dimensionless, so it can be averaged across designs. Note that $P^{*}$ is
the best frontier the agent found on that design, not the true optimum, so
$s = 1$ does not imply optimality. A resource type is not scored for a design if
the design never uses it (e.g., no DSPs) or if no point improves on the
baseline ($\mathrm{HV}(P^{*}) = \mathrm{HV}(\{b\})$).

\paragraph{Per-iteration and cumulative scores.}
We use two ways to choose $P$ in Eq.~\eqref{eq:score}. The
\emph{per-iteration} score $s_i = s(S_i)$ evaluates the frontier proposed by
iteration $i$ alone; it can decrease if a later iteration proposes a worse set
of points. The \emph{cumulative} score $s^{\mathrm{cum}}_i =
s\bigl(\bigcup_{j \le i} S_j\bigr)$ evaluates the best frontier found so far.
Because the pool only grows, $s^{\mathrm{cum}}_i$ is non-decreasing and reaches
exactly $1$ at $i = T$.

\paragraph{Scores against agent cost and time.}
To compare designs with different iteration costs, we also index the score by a
budget $\beta$, either cumulative agent cost (\$) or cumulative agent wall-clock
time (s). If iteration $i$ finishes at cumulative spend $c_i$, the score of a
design at budget $\beta$ is the score of the last iteration with $c_i \le \beta$,
i.e., $s^{\mathrm{cum}}_{i}$ or $s_{i}$ depending on the variant, held constant
until the next iteration completes; the design has no value before its first
iteration finishes. We evaluate every design on a common uniform grid of
budgets, and report the mean across the designs of a benchmark suite at each
grid point, taken over the designs that have a value there. Since the set of
contributing designs changes with $\beta$, the mean is not itself monotone even
when each design's cumulative score is.

\paragraph{Baseline coverage.}
For the per-iteration score we additionally flag iterations whose front
contains no point that is at least as good as the baseline in both latency and
the resource under consideration. Such an iteration can obtain a positive
score by improving one objective while regressing on the other, and the flag
identifies these cases.
